\section{Thiết kế Cloud Control Plane}
\label{sec:cloud-design}

Cloud Control Plane gồm các application service xử lý yêu cầu đồng bộ từ người
dùng và các handler xử lý thông điệp bất đồng bộ từ Edge
(Hình~\ref{fig:cloud-components}).

\begin{figure}[H]
\centering
\resizebox{\textwidth}{!}{%
\begin{tikzpicture}[font=\small,
  b/.style={draw, rounded corners, align=center, minimum height=0.9cm, text width=3.3cm, fill=blue!5},
  h/.style={draw, rounded corners, align=center, minimum height=1.0cm, text width=3.3cm, fill=orange!8},
  r/.style={draw, align=center, minimum height=0.9cm, fill=gray!8},
  lbl/.style={font=\scriptsize, fill=white, inner sep=1.5pt}]
  % Tầng vào
  \node[b, fill=green!6, text width=16cm] (entry) at (6,3.2) {REST Router \qquad | \qquad WebSocket Router \qquad | \qquad MCP Tools};
  % Application service
  \node[b] (fleet) at (0,1.6) {FleetService};
  \node[b] (deploy) at (4,1.6) {DeploymentService};
  \node[b] (monitor) at (8,1.6) {MonitoringService};
  \node[b] (realtime) at (12,1.6) {RealtimeService};
  % MQTT
  \node[r, fill=purple!6, text width=15.6cm] (mqtt) at (6,-0.2) {MQTT Broker};
  % Handler
  \node[h] (state) at (0,-2.0) {NodeStateHandler};
  \node[h] (devt) at (4,-2.0) {DeploymentEvent\\Handler};
  \node[h] (tele) at (8,-2.0) {TelemetryHandler};
  \node[h] (evt) at (12,-2.0) {EdgeEventHandler};
  \node[h, fill=red!6, text width=3.7cm] (live) at (-4.2,-2.0) {NodeLiveness\\Service (quét 10 s)};
  % Lưu trữ
  \node[r, text width=10cm] (pg) at (2,-4.0) {PostgreSQL: clusters, nodes, runtimes, deployments, node\_logs};
  \node[r, text width=3.3cm] (redis) at (10,-4.0) {Redis: telemetry};

  \foreach \n in {fleet, deploy, monitor, realtime} \draw[flowarrow] (entry.south -| \n) -- (\n);
  \draw[flowarrow] (deploy) -- node[lbl, right]{CommandGateway: publish lệnh} (deploy |- mqtt.north);
  \draw[flowarrow, <-] (realtime) -- node[lbl, right]{edge/\#} (realtime |- mqtt.north);
  \foreach \n in {state, devt, tele, evt} \draw[flowarrow] (\n |- mqtt.south) -- (\n);
  \draw[flowarrow, dashed] (state) -- (state |- pg.north);
  \draw[flowarrow, dashed] (devt) -- (devt |- pg.north);
  \draw[flowarrow, dashed] (evt.south) -- ++(0,-0.5) -| ([xshift=2.5cm]pg.north);
  \draw[flowarrow, dashed] (tele) -- (tele |- redis.north);
  \draw[flowarrow, dashed] (live) |- (pg.west);
  \draw[flowarrow, dashed] (fleet.west) -| ([xshift=-0.3cm]live.north);
\end{tikzpicture}}
\caption{Các thành phần của Cloud Control Plane}
\label{fig:cloud-components}
\end{figure}

\subsection{Fleet Service}

FleetService hiện thực các use case trên phân cấp Cluster $\rightarrow$ Node
$\rightarrow$ Runtime: tạo cluster (từ chối tên trùng), đăng ký node vào cluster
(từ chối tên trùng trong cùng cluster), liệt kê và tra cứu. Node được đăng ký
trước trên Cloud; khi Agent với \texttt{node\_id} tương ứng bắt đầu gửi
NodeState, Cloud cập nhật thông tin phần cứng. Bản tin NodeState từ node chưa
đăng ký bị ghi log và bỏ qua, tránh việc thiết bị lạ tự chen vào hệ thống.

\subsection{Deployment Service}

DeploymentService ghi desired state rồi phát lệnh; kết quả được Edge báo về bất
đồng bộ. Do đó các API tạo, cập nhật, gỡ bỏ trả mã HTTP \texttt{202 Accepted}.
Các quy tắc chính:

\begin{itemize}
  \item \textbf{Tạo mới:} kiểm tra node tồn tại. Nếu đã có runtime cùng tên trên
        node với deployment còn hiệu lực: spec giống hệt thì trả lại deployment cũ
        (gửi lại idempotent); spec khác thì báo xung đột và yêu cầu dùng PUT.
        Ngược lại, sinh \texttt{runtime\_id}, tạo bản ghi Runtime (CREATING) và
        Deployment, rồi phát lệnh \texttt{deploy}.
  \item \textbf{Cập nhật:} chỉ ghi đè những trường được gửi lên, đồng bộ bản ghi
        Runtime và phát lệnh \texttt{update}.
  \item \textbf{Gỡ bỏ:} chuyển REMOVING và phát lệnh \texttt{remove}; bản ghi
        Deployment được giữ làm lịch sử, bản ghi Runtime bị xóa khi Edge xác nhận.
  \item \textbf{Điều khiển:} start/stop/restart chỉ áp dụng cho runtime có
        deployment còn hiệu lực.
\end{itemize}

Trước khi publish, service luôn gọi \texttt{dispatched()} và lưu deployment, nên
trạng thái IN\_PROGRESS và \texttt{last\_command\_id} đã bền vững trước khi lệnh
rời khỏi Cloud. CommandGateway đóng gói việc publish lên topic
\path{edge/{cluster}/{node}/command/{type}}.

\subsection{Các handler bất đồng bộ}

\begin{itemize}
  \item \textbf{NodeStateHandler} tiêu thụ \path{edge/+/+/state}: cập nhật
        heartbeat, trạng thái Agent, thông tin phần cứng và đồng bộ trạng thái,
        \texttt{container\_id} của từng runtime do node báo cáo.
  \item \textbf{DeploymentEventHandler} tiêu thụ \path{event/deployment}: tìm
        deployment theo \texttt{deployment\_id} hoặc \texttt{runtime\_id}, áp dụng
        kết quả qua \texttt{apply\_result}, cập nhật trạng thái runtime hoặc xóa
        runtime khi đã REMOVED.
  \item \textbf{TelemetryHandler} tiêu thụ \path{telemetry/{kind}}: kiểm tra schema
        theo loại, lưu mẫu vào Redis; mẫu runtime được khóa theo
        \texttt{runtime:\{runtime\_id\}}.
  \item \textbf{EdgeEventHandler} tiêu thụ \path{event/+}: ghi mọi sự kiện vào nhật
        ký của node (\texttt{node\_logs}).
  \item \textbf{NodeLivenessService} quét định kỳ (10~s) và chuyển các node có
        heartbeat cũ hơn \texttt{node\_heartbeat\_timeout\_sec} (30~s) sang OFFLINE.
\end{itemize}

\subsection{Monitoring và Realtime Service}

MonitoringService cung cấp mẫu telemetry mới nhất của một node (theo từng loại),
lịch sử gần đây của một loại và danh sách event. RedisTelemetryRepository lưu
\texttt{telemetry:\{node\}:latest} dạng hash và \texttt{telemetry:\{node\}:\{kind\}}
dạng list giới hạn 720 mẫu (tương đương 1 giờ với chu kỳ 5~s), có TTL 1 giờ.

RealtimeService đăng ký \path{edge/#} trên broker và phân phối thông điệp tới các
WebSocket client theo bộ lọc topic của từng client (hỗ trợ wildcard MQTT
\texttt{+} và \texttt{\#}). Client có thể đổi bộ lọc khi đang kết nối bằng thông
điệp \texttt{subscribe}/\texttt{unsubscribe}; xác nhận \texttt{ack} đi qua cùng
hàng đợi để giữ thứ tự với dữ liệu.

\subsection{Triển khai tiến trình}

Cloud chạy như một tiến trình duy nhất trên cổng 8000, phục vụ REST, WebSocket,
MCP và embedded worker (\texttt{CLOUD\_EMBEDDED\_WORKER=true}). Khi cần mở rộng,
worker có thể tách thành tiến trình riêng (\texttt{apps/cloud/worker}) trong khi API
được nhân bản. Chế độ \texttt{CLOUD\_STORAGE\_BACKEND=memory} thay PostgreSQL và
Redis bằng repository trong bộ nhớ, phục vụ phát triển và kiểm thử tự động.
